\section{Nietzsche and the Gospel of Buddhism}

\begin{quotex}
The student will, indeed, find that nearly every thought expressed in Buddhism and Hindu literature finds expression in the Western world also; and it could not be otherwise for the value of these thoughts is universal, and therefore they could not be more Oriental than Western; the East has advanced beyond the West only in their wider and fuller acceptance.
\flright{\textsc{Ananda Coomaraswamy}}
\end{quotex}


Years ago, before I had known much about Rene Guenon and Tradition, I read Coomaraswamy's book \textit{Buddha and the Gospel of Buddhism}. Published it 1916, it was before Coomaraswamy's encounter with Guenon. Nevertheless, it foreshadowed a Traditional outlook, since Coomaraswamy took pains to point out the correspondences between Buddhism and both Hinduism and Christian mysticism. Oddly enough, and it struck me so at the time, Coomaraswamy also related it to the ideas of the poets \textbf{Friedrich Nietzsche}, \textbf{Walt Whitman}, and \textbf{William Blake}. Since they were three of my favorite poets, I did consider that a plus.

While the author of the introduction (written much later) is perplexed at the mention of Nietzsche, it would be better to explore what Coomaraswamy sees in him, whether he sees something deeper in Nietzsche or is simply off track. That author also objects to Coomaraswamy's use of the word ``Aryan'', although it was perfectly natural at the time the book was written. It does point out how arbitrary prejudices prevent us from seeing something purely. For example, I used to be prejudiced against the middle ages. A Jewish woman, who sometimes led the meditations in our Buddhism study group, had a PhD in Medieval History and convinced me it was much more that a long period of intellectual darkness.

A few years later Coomaraswamy published an essay on the ``Cosmopolitan View of Nietzsche'', in which he explains the religion of Idealistic Individualism, which he calls the religion of modern Europe. He sees Nietzsche as a poet, not a philosopher, and a mystic, speaking in the same voice as Blake and Whitman. Coomaraswamy explains his hermeneutic principle:

\begin{quotex} Nietzsche's originality does not consist in an incomprehensible and unnatural novelty, but in a poetic restatement of a very old position.
\end{quotex}


Given that, Coomaraswamy can ignore the more simplistic interpretations of his thought. Nietzsche was a prophet with a ``passionate protest against unworthy values, Pharisaic virtue, and snobism''. He then moves on to this:

\begin{quotex}
Of special significance is the beautiful doctrine of the Superman, so like the Chinese concept of the Superior Man, and the Indian maha purush, bodhisattva, and jivan-mukta.
\end{quotex}


By ignoring Nietzsche's false philosophical positions, Coomarswamy finds an impulse that is more ancient. ``The Superman, whose virtue stands beyond good and evil, is at once the flower and the leader and saviour of men.'' His ``actions are no longer good or bad, but proceed from his freed nature.'' Nietzsche's objection to the reward of Heaven and Hell can be related to Guenon's idea of deliverance as superior to salvation. Coomaraswamy then deals with the doctrine of the Will to Power.

\begin{quotex}
The will to Power asserts that our life is not to be swayed by motives of pleasure or pain, the pairs of opposites, but is to be directed towards its goal, and that goal is the freedom and spontaneity of the jivan mukta. And this is beyond good and evil. This also set out in the \emph{Bhagavad Gita}: the hero must be superior to pity; resolute for the fray, but unattached to the result.
\end{quotex}


Nietzsche's teaching is \emph{nishkama dharma}\footnote{\url{https://en.wikipedia.org/wiki/Nishkam\_Karma}}, that is, selfless and desireless action, without expectation of result. The will to power is not ``to yield to all the promptings of the senses, to be the slave of caprice.'' There can be no altruism if there is nothing external to myself. ``Physician, heal thyself'' means the highest duty is self-realisation. Or to quote Jesus, ``first cast out the mote from thine own eye.'' Coomaraswamy continues:

\begin{quotex}
The inferior man regulates his life by externals: inasmuch as he is constrained by desire for long life, reputation, riches, rank or offspring, he is not free. The superior man is of another sort, and of him it may be said, with Chuan Tzu, ``that they live in accordance with their own nature. In the whole world they have no equal. They regulate their life by inward things.''
\end{quotex}


So the inferior man believes all his actions should be praiseworthy. On the contrary:

\begin{quotex}
The pattern of man's behaviour is not to be found in any code, but in the principles of the universe, which is continually revealing to us \emph{its own nature}... There exists a voluptuousness that is not sensuality, a passion for power that is not self-assertion, and a selfishness that is more generous than any altruism.
\end{quotex}


Ascetism is for the strong, who choose it. ``Life should grow colder towards the summit.'' That is why ``a Brahman should do nothing for the sake of enjoyment.'' Coomaraswamy concludes with this interesting claim:

\begin{quotex}
Those who have comprehended the decline and fall of Western civilization will recognize in Nietzsche the reawakening of the conscience of Europe.
\end{quotex}


Rather than end here, I want to connect a thought from his essay ``Young India''.

\begin{quotex}
We stand in the West at the close of the great cycle of Christian civilization which attained its zenith, let us say, in the twelfth or thirteenth century, when the creative will of man swept far beyond its personal boundaries, striving to establish an order in the outer world to correspond with the universal order to he world of imagination or eternity.
\end{quotex}


Coomaraswamy clearly holds that Medieval Europe was the peak, or at least the last peak, of Western civilization. Can we really see Nietzsche as the prophet to reawaken the conscience of Europe? Specifically, we must first, as he claims, comprehend the decline and fall Western civilization which had reached its peak in the thirteenth century. This is how he understands that decline:

\begin{quotex}
From the thirteenth to the twentieth century one can follow the progressive decay of life: the ever fainter expression of the creative will, loosening social integration, the substitution of contract for status, the advancement of material and moral to the exclusion of spiritual values, the decline of vision, up to this present hour of pure chaos, when life and art are evidence of centuries of aimlessness.
\end{quotex}



\flrightit{Posted on 2012-09-21 by Cologero}

\begin{center}* * *\end{center}

\begin{footnotesize}
\begin{sffamily}
\texttt{Audrey on 2012-09-21 at 15:37 said:}

Interesting opening quote, although I disagree with it..the East and West serve differently, and whereas the overarching union is one, that is practically not usable unless you are conversing among Buddhas. Likewise the naive notion that the East is more mature, or embraced things more fully.. And if you permit me, the romanticized notion of a distant path of Europe. EU nobility still rules the US behind curtains and veils. A peak is a peak, and a valley is a valley, regardless of humans and their pefception of it, ie time.

\hfill

\texttt{Audrey on 2012-09-21 at 15:55 said:}

Also, Whitman and Blake were influenced by a wave of Occidentalism which fashionably hit the West, so these so called points of contact ard largdly fabricated by man. I actually just read a few hings by Oshawa which even in English ar decidedly, mentally foreignto the core, where a human notpart of that core will beunable to drink from it, erudite and genuine as the article may be.

\hfill

\end{sffamily}
\end{footnotesize}

